\section{Giới thiệu}

\begin{frame}
  \frametitle{Tín hiệu và phân tích tần số}

  \begin{block}{Tín hiệu}
    Dữ liệu trong khoa học và kỹ thuật thường xuất hiện dưới dạng tín hiệu:
    âm thanh, hình ảnh, sóng điện từ hoặc dữ liệu cảm biến.
  \end{block}

  \onslide<+->
  \begin{itemize}[<+->]
    \item Tín hiệu có thể thay đổi theo thời gian.
    \item Tín hiệu thường gồm nhiều thành phần khác nhau.
    \item Để hiểu và phân tích tín hiệu, người ta cần xác định những thành
    phần tạo nên tín hiệu đó, đặc biệt là các thành phần tần số
  \end{itemize}
\end{frame}

\begin{frame}
  \frametitle{Biến đổi Fourier}

  \begin{block}{Khái niệm}
    Biến đổi Fourier (FT) là phép biến đổi tín hiệu từ miền thời gian sang
    miền tần số.
  \end{block}

  \onslide<+->
  \begin{itemize}[<+->]
    \item Miền thời gian: tín hiệu thay đổi theo thời gian.
    \item Miền tần số: tín hiệu thay đổi theo tần số.
    \item Ví dụ: một bài hát gồm nhiều âm thanh ở các tần số khác nhau.
  \end{itemize}
\end{frame}

\begin{frame}
  \frametitle{Biến đổi Fourier}

  \begin{block}{Công thức}
    Với tín hiệu liên tục $x(t)$, FT được định nghĩa bởi:
    \begin{equation}
      X(f) = \int_{-\infty}^{\infty} x(t)e^{-i2\pi ft}\,dt.
    \end{equation}
  \end{block}
\end{frame}


\begin{frame}
  \frametitle{Biến đổi Fourier rời rạc}

  \begin{block}{Khái niệm}
    Biến đổi Fourier rời rạc (DFT) là phép biến đổi tín hiệu thời gian rời rạc
    sang miền tần số.
  \end{block}
\end{frame}

\begin{frame}
  \frametitle{Biến đổi Fourier rời rạc}
  \begin{block}{Công thức}
    Nếu tín hiệu có $N$ mẫu, DFT được định nghĩa bởi:
    \begin{equation}
      X(k)=\sum_{n=0}^{N-1}x(n)e^{-i\frac{2\pi}{N}kn},
      \quad k=0,1,\ldots,N-1.
    \end{equation}
  \end{block}
\end{frame}

\begin{frame}
  \frametitle{Biến đổi Fourier nhanh}

  \begin{block}{Khái niệm}
    FFT là thuật toán tính DFT nhanh bằng cách chia bài toán lớn thành
    nhiều bài toán nhỏ hơn.
  \end{block}

  \onslide<+->
  \begin{itemize}[<+->]
    \item FFT không làm thay đổi kết quả của DFT.
    \item FFT chỉ tối ưu cách tính toán.
  \end{itemize}
\end{frame}

\begin{frame}
  \frametitle{Biến đổi Fourier nhanh}

  \begin{block}{So sánh độ phức tạp}
    FFT giúp giảm đáng kể số phép toán so với cách tính DFT trực tiếp.
  \end{block}

  \onslide<+->
  \begin{itemize}[<+->]
    \item DFT trực tiếp:
      \begin{equation*}
        O(N^2).
      \end{equation*}

    \item FFT:
      \begin{equation*}
        O(N\log N).
      \end{equation*}

    \item Với dữ liệu lớn, FFT nhanh hơn rất nhiều.
  \end{itemize}
\end{frame}

\begin{frame}
  \frametitle{Ứng dụng của FFT}

  FFT được dùng để phân tích nhiều loại tín hiệu:

  \onslide<+->
  \begin{itemize}[<+->]
    \item \alert<.>{Y học}: tín hiệu tim ECG, tín hiệu não EEG.
    \item \alert<.>{Kỹ thuật}: tín hiệu cảm biến.
    \item \alert<.>{Viễn thông}: tín hiệu sóng vô tuyến, tín hiệu mạng.
    \item \alert<.>{Xử lý đa phương tiện}: âm thanh, hình ảnh, video.
  \end{itemize}
\end{frame}

\begin{frame}
  \frametitle{Biến đổi Fourier ngược}
  \begin{block}{Khái niệm}
    Biến đổi Fourier ngược (IFT) là phép biến đổi tín hiệu trong miền tần số
    thành tín hiệu trong miền thời gian.
  \end{block}
\end{frame}

\begin{frame}
  \frametitle{Biến đổi Fourier ngược}
  \begin{block}{Công thức}
    Với tín hiệu liên tục $X(f)$, IFT được định nghĩa bởi:
    \begin{equation}
      x(t)=\int_{-\infty}^{\infty}X(f)e^{i2\pi ft}\,df.
    \end{equation}
  \end{block}
\end{frame}

\begin{frame}
  \frametitle{Biến đổi Fourier rời rạc ngược}

  \begin{block}{Khái niệm}
    Biến đổi Fourier rời rạc ngược (IDFT) là phép biến đổi tín hiệu miền tần
    số rời rạc sang miền thời gian.
  \end{block}
\end{frame}

\begin{frame}
  \frametitle{Biến đổi Fourier rời rạc ngược}
  \begin{block}{Công thức}
    Nếu tín hiệu có $N$ mẫu, IDFT được định nghĩa bời:

    \begin{equation}
      x(n)=\frac{1}{N}\sum_{k=0}^{N-1}X(k)e^{i\frac{2\pi}{N}kn},
      \quad n=0,1,\ldots,N-1.
    \end{equation}
  \end{block}
\end{frame}

\begin{frame}
  \frametitle{Lịch sử phát triển}

  \begin{block}{Các cột mốc chính}
    \begin{itemize}[<+->]
      \item \alert<.>{1805}: C. F. Gauss sử dụng một phương pháp tương tự FFT
        khi tính quỹ đạo các tiểu hành tinh \cite{heideman1984gauss}.

      \item \alert<.>{1942}: Danielson và Lanczos đề xuất phương pháp phân rã
        DFT trong tinh thể học tia X \cite{danielson1942}.

      \item \alert<.>{1965}: Cooley và Tukey công bố thuật toán FFT hiện đại
        \cite{cooley1965algorithm}.
    \end{itemize}
  \end{block}
\end{frame}

\begin{frame}
  \frametitle{Tầm quan trọng}

  \begin{block}{Nhận xét}
    FFT là một trong những thuật toán quan trọng nhất trong khoa học máy tính
    và xử lý tín hiệu số.
  \end{block}

  \onslide<+->
  \begin{itemize}[<+->]
    \item Giảm độ phức tạp từ $O(N^2)$ xuống $O(N\log N)$.
    \item Được ứng dụng rộng rãi trong nhiều lĩnh vực.
    \item Được xem là một trong 10 thuật toán tiêu biểu của thế kỷ XX
      \cite{10.1109/MCISE.2000.814652}.
  \end{itemize}
\end{frame}
